\begin{lemma}
For any fixed red row map $\rho: \operatorname{Fin}n \to \operatorname{Fin}m$, target set $S \subseteq \operatorname{Fin}m$, and set of original vertices $T \subseteq \operatorname{Fin}n$, let
\[
X_{\rm all}
=\{\phi:\operatorname{Fin}n\hookrightarrow
\operatorname{Fin}m\times\operatorname{Fin}m:
(\phi(v))_1=\rho(v)\text{ for every }v\}
\]
and let $X_{\rm good} \subseteq X_{\rm all}$ be the subset for which $(\phi(v))_2 \in S$ for every $v \in T$.
For each row $i\in\operatorname{Fin}m$, put
\[
R_i=\{v:\rho(v)=i\},\qquad q_i=|R_i|,\qquad
T_i=T\cap R_i,\qquad h_i=|T_i|.
\]
Then the cardinality of $X_{\rm good}$ is given by the product over rows of the favorable factors:
\[
|X_{\rm good}| = \prod_{i\in\operatorname{Fin}m} (|S|)_{h_i} (m-h_i)_{q_i-h_i}.
\]
\end{lemma}
\begin{proof}
For each row \(i\), let \(\mathcal G_i\) be the set of injective maps
\[
b_i:R_i\hookrightarrow \operatorname{Fin}m
\]
such that \(b_i(v)\in S\) for every \(v\in T_i\).  We first identify
\(X_{\rm good}\) with the product of these rowwise sets.  Given
\(\phi\in X_{\rm good}\), define \(b_i(v)=(\phi(v))_2\) for \(v\in R_i\).
The same argument as for the unrestricted row decomposition shows that
\(b_i\) is injective: if two vertices in the same row have the same blue
coordinate, then their full pairs under \(\phi\) are equal, so the vertices
are equal.  Since \(\phi\in X_{\rm good}\), every \(v\in T_i\) has
\((\phi(v))_2\in S\), so \(b_i\in\mathcal G_i\).

Conversely, given a family \((b_i)_i\) with \(b_i\in\mathcal G_i\), assemble
\[
\phi(v)=(\rho(v),b_{\rho(v)}(v)).
\]
This is well-defined because \(v\in R_{\rho(v)}\).  It has first coordinate
\(\rho(v)\) for every \(v\), so it lies in \(X_{\rm all}\).  It is globally
injective: if \(\phi(v)=\phi(w)\), equality of first coordinates puts both
vertices in the same row \(i\), and equality of second coordinates plus
injectivity of \(b_i\) gives \(v=w\).  Finally, if \(v\in T\), then
\(v\in T_{\rho(v)}\), so \(b_{\rho(v)}(v)\in S\); hence the assembled
\(\phi\) lies in \(X_{\rm good}\).  These two constructions are inverse
pointwise.  Therefore
\[
|X_{\rm good}|=\prod_{i\in\operatorname{Fin}m}|\mathcal G_i|.
\]

It remains to compute \(|\mathcal G_i|\) for a fixed row \(i\).  If
\(h_i>|S|\), then the \(h_i\) vertices of \(T_i\) would have to receive
distinct blue coordinates inside the \(|S|\)-element set \(S\).  This is
impossible by injectivity, so \(|\mathcal G_i|=0\).  The factor
\((|S|)_{h_i}\) is also \(0\), so the claimed row formula holds in this
case.

Next suppose \(q_i>m\).  Then no injective map from the \(q_i\)-element set
\(R_i\) into \(\operatorname{Fin}m\) exists, so \(|\mathcal G_i|=0\).  On the
right hand side, if \(h_i>|S|\), the previous paragraph already gives
\((|S|)_{h_i}=0\).  Otherwise \(h_i\le |S|\), and since \(h_i\le q_i\) and
\(q_i>m\), we have \(q_i-h_i>m-h_i\).  Thus
\((m-h_i)_{q_i-h_i}=0\).  In all subcases the product
\((|S|)_{h_i}(m-h_i)_{q_i-h_i}\) is \(0\), matching
\(|\mathcal G_i|\).

We are left with the ordinary case \(h_i\le |S|\) and \(q_i\le m\).  First
choose the blue coordinates of the \(h_i\) constrained vertices in \(T_i\).
They must be distinct elements of \(S\), so there are \((|S|)_{h_i}\)
choices.  After these coordinates are used, the remaining \(q_i-h_i\)
vertices in \(R_i\setminus T_i\) may use any blue coordinates not already
used; there are \(m-h_i\) available coordinates and the choices must be
distinct, giving \((m-h_i)_{q_i-h_i}\) choices.  Multiplying these independent
stages gives
\[
|\mathcal G_i|=(|S|)_{h_i}(m-h_i)_{q_i-h_i}.
\]
Combining this row formula for every \(i\) with the product decomposition of
\(X_{\rm good}\) proves
\[
|X_{\rm good}|
=\prod_{i\in\operatorname{Fin}m} (|S|)_{h_i}(m-h_i)_{q_i-h_i}.
\]
\end{proof}
